% Plantoir house style: a handout, worksheet, review sheet or exercise.
% ONE source, up to three outputs:
%   this file                  -> Handout.pdf            the questions
%   "Handout - Answers.tex"    -> Handout - Answers.pdf  final answers (defines \showanswers)
%   "Handout - Hints.tex"      -> Handout - Hints.pdf    the hints alone (defines \showhints)
% Copy the files you need to where the teacher keeps handouts, rename them
% to the handout's name, and change the \input line in each wrapper to match.
% Fill in the lines marked "Fill in" below. Delete the hints block (and the
% hints wrapper) when the exercise has no hints.
%
% This starter is a WORKING EXAMPLE of a finished sheet, not an empty
% template: worked examples in peach boxes with substituted values in red
% and side notes beside the steps, a table and a graph inside a box, a
% table with negatives in an N column, a graph drawn to the grid's edge
% with arrowheads, a context graph that stops at its end points, fill-in
% blanks, and the answer key on the SAME tables and grids. Keep those forms
% and change the mathematics. The skill's sections 3 to 5 and 9 to 13 say why.
\documentclass[11pt]{article}
\usepackage[letterpaper, left=0.75in, right=0.75in, top=1.05in, bottom=0.9in,
            headheight=14pt, headsep=0.3in]{geometry}
\usepackage{lmodern}
\usepackage[T1]{fontenc}
\usepackage{amsmath, amssymb}
\usepackage{xcolor}
\usepackage{enumitem}
\usepackage{tasks}
\usepackage{fancyhdr}
\usepackage{lastpage}
\usepackage{pgfplots}
\usepackage{array}
\usepackage{collcell}     % the C column's minus-sign warning (below)
\usepackage{needspace}
\usepackage[most]{tcolorbox}
\usepackage{nicematrix}   % tables: every rule drawn after the heading's shade (the skill's section 9)
\pgfplotsset{compat=1.18}

% ---------- Fill in ----------
% The school's name comes ONLY from print_school_name in the course's
% settings, or from the teacher in this conversation. Never a placeholder:
% with no name, leave the braces empty and the corner prints nothing.
\newcommand{\schoolname}{}
\newcommand{\coursecode}{COURSE}        % the course's code, in capitals
\newcommand{\handouttitle}{Handout Title}

% ---------- Questions, answers or hints ----------
\newif\ifanswers
\ifdefined\showanswers \answerstrue \fi
\newcommand{\qa}[2]{\ifanswers #2\else #1\fi}
\newcommand{\onlyq}[1]{\ifanswers\else #1\fi}
\newcommand{\onlya}[1]{\ifanswers #1\fi}
% The hints sheet (the skill's section 13): the hints alone, one per
% question, in the block right after \begin{document}. Nothing else in the
% file changes for hints, so the questions and the answers compile exactly
% as they would without them.
\newif\ifhints
\ifdefined\showhints \hintstrue \fi
\newcommand{\hint}[2]{\item[Question #1] #2}

% A blank students write on, in a question (the skill's section 4): an
% invisible strut makes its line about two lines tall, so a wrapped prompt
% still leaves a full line of writing room above the blank. \fillin{answer}
% or \fillin[5cm]{answer}: the student copy shows the blank, the answer key
% writes the answer on the same line, at the same height.
\newcommand{\fillin}[2][3.2cm]{\ \qa{\rule[-0.2em]{#1}{0.4pt}}{\makebox[#1]{#2}\llap{\rule[-0.2em]{#1}{0.4pt}}}\rule[-0.2em]{0pt}{2.1em}}

\ifanswers
  \newcommand{\doctitle}{\handouttitle\ --- Answers}
\else\ifhints
  \newcommand{\doctitle}{\handouttitle\ --- Hints}
\else
  \newcommand{\doctitle}{\handouttitle}
\fi\fi

% ---------- Header and footer (matches pages printed from the class website) ----------
\pagestyle{fancy}
\fancyhf{}
\renewcommand{\headrulewidth}{0pt}
\newcommand{\fillblank}[2]{#1 \rule[-0.15em]{#2}{0.4pt}}
% Lay the corners out as the course's print settings do (the skill's
% sections 2 and 5). The blanks are print_blanks' list in the order Name,
% Date, Class #; with no list, Name and Date, as here. Add
% \quad\fillblank{Class \#}{1.4cm} when the list has class_number.
\fancyhead[L]{\small\sffamily\color{gray}\schoolname}
\fancyhead[R]{\small\sffamily\color{gray}\fillblank{Name}{3.2cm}\quad\fillblank{Date}{2.2cm}}
\fancyfoot[L]{\small\sffamily\color{gray}\doctitle\ $\cdot$ \coursecode}
\fancyfoot[R]{\small\sffamily\color{gray}Page \thepage\ of \pageref{LastPage}}
\setlength{\parindent}{0pt}

% ---------- Colours and boxes ----------
\definecolor{why}{gray}{0.4}
\definecolor{peach}{RGB}{253,244,233}
\definecolor{peachframe}{RGB}{226,160,96}
\definecolor{exampletitle}{RGB}{176,38,30}
\definecolor{sub}{RGB}{200,30,30}      % a substituted value, in a worked example
\definecolor{note}{RGB}{20,70,140}     % a side note beside a step
\definecolor{tablehead}{RGB}{225,232,245}   % a table's heading row

% Worked example box
\newtcolorbox{example}[1]{
  enhanced, breakable=false,
  colback=peach, colframe=peachframe, boxrule=0.8pt, arc=3mm,
  left=3mm, right=3mm, top=2mm, bottom=2mm,
  before skip=0.9em, after skip=0.6em,
  fontupper=\normalsize,
  title={#1}, coltitle=exampletitle, colbacktitle=peach,
  fonttitle=\large\bfseries\sffamily,
  titlerule=0pt, toptitle=1mm, bottomtitle=-1mm,
}
% Inside a worked example: a step's method as a small heading, a note beside
% a step (in an alignat, after &&), and a substituted value in red.
\newcommand{\method}[1]{\medskip{\bfseries\sffamily #1}\par\smallskip}
\newcommand{\sidenote}[1]{\text{\small\sffamily\color{note}#1}}
\newcommand{\sv}[1]{{\color{sub}#1}}
% A small framed note set beside an example, in a [t] minipage
\newtcolorbox{notebox}{colback=white, colframe=peachframe, boxrule=0.5pt, arc=1.5mm,
  left=2mm, right=2mm, top=1mm, bottom=1mm, fontupper=\small\sffamily}
% A grey box of instructions before a group of questions
\newtcolorbox{groupnote}{colback=gray!8, colframe=gray!45, boxrule=0.6pt, arc=2mm,
  left=3mm, right=3mm, top=1.5mm, bottom=1.5mm, before skip=0.4em, after skip=0.8em,
  fontupper=\small}

% Section heading for the answer key, where the worked examples are hidden
\newcommand{\strand}[1]{\bigskip{\large\bfseries\sffamily #1}\par\smallskip}

% ---------- Tables (the skill's section 9) ----------
% Every value sits in a row \valuepitch tall, centred in it, so the first
% value is half a row below the heading rule and the last half a row above
% the bottom rule. The value columns of a table are one width
% (\setvaluecol{the widest value in either column} sets \valuecol to it plus
% 0.5cm). C is a centred value column, for a column with no negative value
% and not mixing whole numbers with decimals. N{$18$} is for every other
% column of numbers: each value right-aligned in a box as wide as the
% column's widest value WITHOUT its sign, the box centred, so the digits
% share one trailing edge and a minus sign juts out to the left.
% A C cell holding a minus sign prints a small red "use N" beside it, so a
% centred column of negatives shows on the rendered page (the skill's
% section 8). An N column never does. \renewcommand{\valuepitch}{3em} for a
% table with room to write.
\renewcommand{\arraystretch}{1.1}
\newcommand{\valuepitch}{1.5em}
\newcommand{\rowstrut}[1]{\rule[\dimexpr0.32em-0.5\dimexpr#1\relax\relax]{0pt}{#1}}
\newlength{\valuecol}
\newcommand{\setvaluecol}[1]{\settowidth{\valuecol}{#1}\addtolength{\valuecol}{0.5cm}}
\makeatletter
\newcommand{\centredvalue}[1]{%
  \begingroup\edef\x{\endgroup\noexpand\in@{\detokenize{$-}}{\detokenize{#1}}}\x
  \ifin@\else\begingroup\edef\x{\endgroup\noexpand\in@{\detokenize{$ -}}{\detokenize{#1}}}\x\fi
  \ifin@\llap{\tiny\sffamily\color{red}use N\ }\fi#1}
\makeatother
\newcolumntype{C}{>{\rowstrut{\valuepitch}\collectcell\centredvalue}w{c}{\valuecol}<{\endcollectcell}}
\newcolumntype{N}[1]{>{\rowstrut{\valuepitch}\settowidth{\dimen0}{#1}\hbox to\dimen0\bgroup\hss}w{c}{\valuecol}<{\egroup}}
% One heading-row height for tables set side by side
\newcommand{\headstrut}{\rule[-1.1em]{0pt}{2.6em}}

% ---------- Numbered questions ----------
\newenvironment{questions}
  {\begin{list}{}{\setlength{\leftmargin}{2.2em}\setlength{\labelwidth}{1.8em}
                  \setlength{\labelsep}{0.4em}\setlength{\itemsep}{0.7em}}}
  {\end{list}}
% Optional argument: lines to keep together so a question is never split from its parts
\newcommand{\q}[2][5]{\needspace{#1\baselineskip}\item[\textbf{#2.}]}

\settasks{label=\alph*), label-width=1.4em, item-indent=1.8em, after-item-skip=0.35em}

% ---------- Graph styles (the skill's section 12) ----------
% Tick labels are drawn above the plotted lines on a white patch, so a line
% that passes a label never runs through its digits. The axes run a little
% past the grid, so a line's arrowhead at the grid's edge is clear of the
% axis arrows. Leave a tick label out where an arrowhead ends on it
% (yticklabels={,$-4$,...}).
\pgfplotsset{
  layers/labels on top/.define layer set={axis background, axis grid, axis ticks, axis lines,
    main, axis tick labels, axis descriptions, axis foreground}{},
  gridgraph/.style={
    set layers=labels on top,
    axis lines=middle, grid=both,
    major grid style={gray!35}, minor grid style={gray!18},
    minor tick num=1,
    tick label style={font=\scriptsize, fill=white, fill opacity=0.85, text opacity=1, inner sep=1pt},
    xlabel=$x$, ylabel=$y$,
    xlabel style={anchor=west, xshift=0.25cm}, ylabel style={anchor=south, yshift=0.25cm},
    axis line style={-latex, shorten >=-0.25cm},
    clip mode=individual,
  },
  % A graph that models a situation: axes from zero at the bottom left,
  % labelled with what they measure and its units, no arrowheads on the line.
  contextgraph/.style={
    set layers=labels on top,
    axis lines=left, grid=both,
    major grid style={gray!35}, minor grid style={gray!18},
    tick label style={font=\scriptsize},
    label style={font=\small},
    axis line style={-latex, shorten >=-0.25cm},
    clip mode=individual,
  },
  exampleline/.style={blue!70!black, thick, samples=2, latex-latex},
  lineA/.style={blue!70!black, thick, samples=2},
  lineB/.style={red!75!black, thick, samples=2},
  lineC/.style={green!50!black, thick, samples=2},
  lineD/.style={orange!85!black, thick, samples=2},
  % a legend whose keys show a plain line, without the arrowheads the lines carry
  answerlegend/.style={legend pos=outer north east, legend cell align=left,
    legend style={font=\small, draw=gray!50}, legend image post style={arrows=-}},
}

\begin{document}

% ---------- The hints sheet: only in "Handout - Hints.pdf" ----------
% One \hint per question that has one. A hint never gives the answer, not
% even through a pointed first step: it points back to the day's idea, in
% the words the lesson used (the skill's section 13).
\ifhints
{\LARGE\bfseries\sffamily \doctitle}\par
\medskip
{\color{why} Try each question first. Open a hint only when you are stuck.
The answers come out at the start of next class.}
\begin{description}[leftmargin=6.5em, style=nextline, itemsep=0.8em]
  \hint{1}{Look back at the worked example ``Substitute and Evaluate''. Each value goes where its letter was, in brackets.}
  \hint{2}{Look back at ``Solve an Equation''. The side notes name each step in the order it was done.}
  \hint{3}{Look back at ``Graph a Line from a Table of Values''. Each $y$ value comes from its own $x$ value.}
  \hint{4}{Look back at what the lesson called the \emph{slope} and the \emph{$y$-intercept}, then at your table and graph from question 3.}
  \hint{5}{Each second the cyclist goes the same distance. The graph stops where the ride stops.}
\end{description}
\end{document}
\fi

{\LARGE\bfseries\sffamily \doctitle}\par
\medskip
\ifanswers
  {\color{why} Final answers only. Equivalent forms are fine (for example, $(0, 3)$ or ``the $y$-intercept is $3$'').}
\else
  {\color{why} Read each worked example, then try the questions that follow it.
  Show all your work.}
\fi

% =====================================================================
\ifanswers
\strand{Substitute and evaluate}
\else
\begin{example}{Substitute and Evaluate}
Evaluate $2x^2 - 3y$ when $x = -3$ and $y = 4$.
\begin{align*}
  & 2x^2 - 3y \\
  ={} & 2(\sv{-3})^2 - 3(\sv{4}) \\
  ={} & 2(9) - 12 \\
  ={} & 6
\end{align*}
\end{example}
\fi

\begin{questions}
\q{1} \onlyq{Evaluate each expression when $x = -2$ and $y = 5$.}
\begin{tasks}(3)
  \task \qa{$3x + y$}{$-1$}
  \task \qa{$x^2 - y$}{$-1$}
  \task \qa{$4y - 2x$}{$24$}
  \task \qa{$-x + 3y$}{$17$}
  \task \qa{$xy + 7$}{$-3$}
  \task \qa{$\dfrac12x - y$}{$-6$}
\end{tasks}
\end{questions}

% =====================================================================
\ifanswers
\strand{Solve an equation}
\else
\begin{example}{Solve an Equation}
Solve $3(x - 2) + 4 = 13$.
\begin{alignat*}{2}
  3(x - 2) + 4 &= 13 \qquad && \sidenote{Expand the brackets.} \\
  3x - 6 + 4 &= 13 && \sidenote{Collect like terms.} \\
  3x - 2 &= 13 && \sidenote{Add 2 to both sides.} \\
  3x &= 15 && \sidenote{Divide both sides by 3.} \\
  x &= 5 &&
\end{alignat*}
\method{Check}
$3(\sv{5} - 2) + 4 = 3(3) + 4 = 13$, so $x = 5$ is right.
\end{example}
\fi

\begin{questions}
\q[1]{2} \onlyq{Solve. Check each answer.}
\begin{tasks}(3)
  \task \qa{$2(x + 3) = 14$}{$x = 4$}
  \task \qa{$5x - 7 = 3x + 9$}{$x = 8$}
  \task \qa{$4(x - 1) - 2 = 18$}{$x = 6$}
\end{tasks}
\end{questions}

% =====================================================================
% A worked example with a table and its graph side by side INSIDE the box:
% two [t] minipages, the right one starting with \vspace{0pt} so both share
% a top edge. The line has no context, so it runs to the grid's edge with an
% arrowhead at each end; its label sits on a peach patch.
\ifanswers
\strand{Graph a line from a table of values}
\else
\begin{example}{Graph a Line from a Table of Values}
\begin{minipage}[t]{0.56\linewidth}
Graph the line $y = x - 2$.

\smallskip
Choose simple values for $x$. Find each $y$ value: when $x = \sv{0}$,
$y = \sv{0} - 2 = -2$.

\medskip
\begin{center}
{\setvaluecol{$-3$}%
\begin{NiceTabular}{N{$3$}N{$3$}}[hvlines]
\CodeBefore \rowcolor{tablehead}{1} \Body
\multicolumn{1}{c}{$x$} & \multicolumn{1}{c}{$y$} \\
$-1$ & $-3$ \\
$0$ & $-2$ \\
$1$ & $-1$ \\
$2$ & $0$ \\
$3$ & $1$ \\
\end{NiceTabular}}
\end{center}

\medskip
Plot the points. Draw a line through them, to the edge of the grid.
\end{minipage}\hfill
\begin{minipage}[t]{0.4\linewidth}
\vspace{0pt}\centering
\begin{tikzpicture}\begin{axis}[gridgraph, x=0.42cm, y=0.42cm,
    xmin=-3, xmax=5, ymin=-6, ymax=4, xtick={-2,2,4}, ytick={-4,-2,2}]
  \addplot[exampleline, domain=-2.9:4.9]{x - 2};
  \addplot[only marks, mark size=1.6pt, blue!70!black] coordinates {(-1,-3) (0,-2) (1,-1) (2,0) (3,1)};
  \node[anchor=west, font=\small, fill=peach, inner sep=1pt] at (axis cs:2.2,-1.6) {$y = x - 2$};
\end{axis}\end{tikzpicture}
\end{minipage}
\end{example}
\fi

\onlyq{%
\begin{groupnote}
\textbf{\sffamily For questions 3 and 4:} use the table you complete in question 3 to draw the graph.
\end{groupnote}}

\begin{questions}
% The table and the grid are written ONCE, outside the switch: the key
% switches only the cells (\qa{}{...}) and the plotted answer (\onlya{...}),
% so both copies have the same table, the same grid and the same size.
% Both columns hold a negative, so both are N columns (N{$3$} and N{$5$}:
% each one's widest value without its sign), their headings centred.
\q[14]{3} \qa{Complete the table of values for $y = -2x + 3$. Then graph the line.}{$y = -2x + 3$:}

\smallskip
\begin{minipage}[t]{0.35\linewidth}\vspace{0pt}\centering
{\setlength{\valuecol}{1.6cm}%
\begin{NiceTabular}{N{$3$}N{$5$}}[hvlines]
\CodeBefore \rowcolor{tablehead}{1} \Body
\multicolumn{1}{c}{$x$} & \multicolumn{1}{c}{$y$} \\
$-1$ & \qa{}{$5$} \\
$0$ & \qa{}{$3$} \\
$1$ & \qa{}{$1$} \\
$2$ & \qa{}{$-1$} \\
$3$ & \qa{}{$-3$} \\
\end{NiceTabular}}
\end{minipage}\hfill
\begin{minipage}[t]{0.6\linewidth}\vspace{0pt}\centering
% It leaves the grid through y = 6 at x = -1.5 and y = -6 at x = 4.5.
\begin{tikzpicture}\begin{axis}[gridgraph, x=0.42cm, y=0.42cm,
    xmin=-4, xmax=6, ymin=-6, ymax=6, xtick={-4,-2,...,6}, ytick={-4,-2,...,4}]
  \onlya{%
  \addplot[lineA, <->, >=latex, domain=-1.45:4.45]{-2*x + 3};
  \addplot[only marks, mark size=1.6pt, blue!70!black] coordinates {(-1,5) (0,3) (1,1) (2,-1) (3,-3)};}
\end{axis}\end{tikzpicture}
\end{minipage}

\q{4} The line $y = -2x + 3$ crosses the $y$-axis at the point \fillin[2.4cm]{$(0, 3)$}.
Its slope is \fillin[2cm]{$-2$}. Each time $x$ goes up by 1, the value of $y$ \fillin[4cm]{goes down by 2}.
\end{questions}

% =====================================================================
% A graph that models a situation stops at the end points the situation
% gives it, with a dot at each and no arrowheads. The distance column mixes
% whole numbers and decimals, so every value has one decimal place
% (0.0, 9.0, 18.0) and it is an N column; the time column is whole numbers
% only, so it is centred. The sentences say 9 m, never 9.0 m.
\begin{questions}
\q[16]{5} A cyclist starts a ride at the start line and rides at a steady speed of
4.5 metres every second for 4 seconds.

\smallskip
\begin{minipage}[t]{0.35\linewidth}\vspace{0pt}\centering
{\setvaluecol{Distance (m)}%
\begin{NiceTabular}{CN{$18.0$}}[hvlines]
\CodeBefore \rowcolor{tablehead}{1} \Body
Time (s) & \multicolumn{1}{c}{Distance (m)} \\
$0$ & $0.0$ \\
$1$ & $4.5$ \\
$2$ & \qa{}{$9.0$} \\
$3$ & \qa{}{$13.5$} \\
$4$ & \qa{}{$18.0$} \\
\end{NiceTabular}}
\end{minipage}\hfill
\begin{minipage}[t]{0.6\linewidth}\vspace{0pt}\centering
\begin{tikzpicture}\begin{axis}[contextgraph, x=1.1cm, y=0.22cm,
    xmin=0, xmax=5, ymin=0, ymax=20, xtick={0,1,...,5}, ytick={0,4,...,20}, minor y tick num=1,
    xlabel={Time (s)}, ylabel={Distance (m)}]
  \onlya{%
  \addplot[lineA, domain=0:4]{4.5*x};
  \addplot[only marks, mark size=1.6pt, blue!70!black] coordinates {(0,0) (1,4.5) (2,9) (3,13.5) (4,18)};}
\end{axis}\end{tikzpicture}
\end{minipage}

\begin{tasks}(1)
  \task \onlyq{Complete the table.}\onlya{The table above.}
  \task \onlyq{Graph the ride on the grid.}\onlya{The graph above: a line from $(0, 0)$ to $(4, 18)$, with no arrowheads.}
  \task How far has the cyclist gone after 2 seconds? \fillin[2.4cm]{9 m}
  \task After how many seconds has the cyclist gone 13.5 m? \fillin[2.4cm]{3 s}
\end{tasks}
\end{questions}

\end{document}
